% !TeX spellcheck = de_DE
%%
% This is an Overleaf template for presentations
% using the TUM Corporate Desing https://www.tum.de/cd
%
% For further details on how to use the template, take a look at our
% GitLab repository and browse through our test documents
% https://gitlab.lrz.de/latex4ei/tum-templates.
%
% The tumbeamer class is based on the beamer class.
% If you need further customization please consult the beamer class guide
% https://ctan.org/pkg/beamer.
% Additional class options are passed down to the base class.
%
% If you encounter any bugs or undesired behaviour, please raise an issue
% in our GitLab repository
% https://gitlab.lrz.de/latex4ei/tum-templates/issues
% and provide a description and minimal working example of your problem.
%%

%% ================================== base ================================== %%
\documentclass[
	german,            % define the document language (english, german)
	aspectratio=169,    % define the aspect ratio (169, 43)
	% handout=2on1,       % create handout with multiple slides (2on1, 4on1)
	% partpage=false,     % insert page at beginning of parts (true, false)
	% sectionpage=true,   % insert page at beginning of sections (true, false)
]{tumbeamer}

\usepackage{booktabs}
\usepackage{tabularx}
\usepackage{enumerate}
\usepackage{float}

% fvextra must be loaded before csquotes (minted pulls it in later)
\usepackage{fvextra}

%% ================================ language ================================ %%
\usepackage{babel}[ngerman]
\usepackage{csquotes}[autostyle]
\usepackage[useregional]{datetime2}
\usepackage{pmboxdraw}

%% ================================== math ================================== %%
\usepackage{amsmath}
\usepackage{amssymb}
\usepackage{ifthen}
\usepackage{siunitx}

%% ================================ graphics ================================ %%
\usepackage{graphicx}
\usepackage[dvipsnames]{xcolor}
\usepackage{pifont}
\usepackage{url}
\usepackage{hyperref}

\graphicspath{ {./resources/} }
\captionsetup{labelformat=empty}

%% ================================== tikz ================================== %%
\usepackage{pgf}
\usepackage{tikz}
\usepackage{pgfplots}
\usepackage{tikz-timing}
\usepackage{circuitikz}
\usepackage{karnaugh-map}

\usetikzlibrary{arrows, arrows.meta, backgrounds, positioning,
	shapes, shapes.geometric, shapes.misc,
	patterns, patterns.meta,
	matrix, fit, calc,
	decorations.pathmorphing,
	overlay-beamer-styles,
	automata, bending,
	riscvproc}

% requires circuitikz >= 1.1.0
% for distros with older distributions, install TeX Live manually
% instead of using your package manager
% see: https://tug.org/texlive/quickinstall.html
\ctikzset{logic ports=ieee}

%% ============================= code listings ============================= %%
\usepackage[cache=true]{minted}
\usepackage{listings}
\usepackage{lstautogobble}

\usemintedstyle{borland}
\setminted{
	fontsize=\small,
	frame=none,
	breaklines=true,
}

%% ============================ colors & styling ============================ %%
\definecolor{altblue}{HTML}{44C8F5}
\definecolor{altgreen}{HTML}{26E600}
\definecolor{altmagenta}{HTML}{EC008C}
\definecolor{altyellow}{HTML}{FFE600}

\renewcommand{\thempfootnote}{\arabic{mpfootnote}}
\setbeamercolor{footnote}{fg=black}
\setbeamercolor{footnote mark}{fg=black}

%% ============================= custom macros ============================= %%
% https://tex.stackexchange.com/a/7045
\newcommand*\circled[1]{\tikz[baseline=(char.base)]{
		\node[shape=circle,draw,inner sep=2pt, font=\scriptsize] (char) {#1};}}

\newcommand*\colorcirc[1]{\tikz[baseline=(char.base)]{
		\node[shape=circle,fill=#1,inner sep=2pt] {};}}

\newcommand\n[1]{\mkern1mu\overline{\mkern-1mu#1}}

%% ================================ metadata ================================ %%
\author{Niklas Ladurner}
\institute{\theChairName\\\theDepartmentName\\\theUniversityName}

\footline{\insertauthor~|~\insertshorttitle~|~\insertshortdate}

\TUMbeamersetup{
	title page = TUM tower,         % style of the title page
	part page = TUM toc,            % style of part pages
	section page = TUM toc,         % style of section pages
	content page = TUM more space,  % style of normal content pages
	tower scale = 1.0,              % scaling factor of TUM tower (if used)
	headline = TUM threeliner,      % which variation of headline to use
	footline = TUM default,         % which variation of footline to use
	% configure on which pages headlines and footlines should be printed
	headline on = {title page},
	footline on = {every page, title page=false},
}

% available frame styles for title page, part page, and section page:
% TUM default, TUM tower, TUM centered,
% TUM blue default, TUM blue tower, TUM blue centered,
% TUM shaded default, TUM shaded tower, TUM shaded centered,
% TUM flags
%
% additional frame styles for part page and section page:
% TUM toc
%
% available frame styles for content pages:
% TUM default, TUM more space
%
% available headline options:
% TUM empty, TUM oneliner, TUM twoliner, TUM threeliner, TUM logothreeliner
%
% available footline options:
% TUM empty, TUM default, TUM infoline


\title{Übung 01: Zahlensysteme}
\subtitle{Einführung in die Rechnerarchitektur}
\date{\DTMdisplaydate{2026}{10}{16}{-1}}

\begin{document}

\maketitle

\begin{frame}[c]{}{}
  \begin{center}
    \LARGE  Keine Garantie für die Richtigkeit der Tutorfolien.

    \Large Bei Unklarheiten/Unstimmigkeiten haben VL/ZÜ-Folien recht!
  \end{center}
\end{frame}

\begin{frame}[c]{Organisatorisches}{}
	\begin{columns}[c]
		\begin{column}{0.6\textwidth}
			\begin{itemize}
				\item Wer bin ich?
				\item Kommunikation:
				\begin{itemize}
					\item Zulip: ERA-Streams vs. Tutoriumsstream vs. DM
					\item E-Mail
					\item Fax (??)
				\end{itemize}
				\item Mitschriften/Folien auf meiner Homepage
				\item Tutoriumszeiten
				\item Artemis-Hausaufgaben: 100\% bedeuten 100\%, keine \emph{hidden tests} nach der Deadline
				\item Laptops mitbringen!
			\end{itemize}
				\end{column}
		\begin{column}{0.4\textwidth}
			\begin{center}
				\LARGE  \href{https://home.in.tum.de/~ladu/}{home.cit.tum.de/\string~ladu/}\\
				\includegraphics[height=0.5\textheight]{website_qr.pdf}
			\end{center}
		\end{column}
	\end{columns}

\end{frame}

\begin{frame}[c]{Motivation}{}
  \begin{itemize}
    \item Rechner arbeiten mit Zahlen: Bilder, Strings, ...
    \item Dezimalsystem ungeeignet: zu viele Zustände (Ziffern)
    \item Strom an und aus mittels Transistoren
    \item zwei Zustände, d.h. Binärsystem
  \end{itemize}
\end{frame}

\begin{frame}[c]{Definition}{}
	\begin{itemize}
		\item Zahlensystem: Bestimmt die Darstellung einer Zahl zu einer gegebenen Basis $B$
		\item Gleiche \emph{Werte} können in unterschiedlichen Systemen unterschiedliche \emph{Darstellungen} haben, bspw. $12_{10}=1100_2=C_{16}$
	\end{itemize}
	
	\begin{equation}
		W=\sum_{i=0}^{n-1}a_i\cdot B^i
	\end{equation}
	
	\hspace{1cm}
	
	\begin{block}{Beispiele}
		\begin{enumerate}
			\item $(1010)_2 = 0\cdot 2^0+1\cdot 2^1+0\cdot 2^2+1\cdot 2^3=2+8=10$
			\item $(763)_8 = 3\cdot 8^0+6\cdot 8^1+7\cdot 8^2=3+48+448=499$
			\item $(123)_{16} = 3\cdot 16^0+2\cdot 16^1+1\cdot 16^2=3+32+256=291$
		\end{enumerate}
	\end{block}
\end{frame}

\begin{frame}[c]{Umwandlungen}{}
	\begin{itemize}
		\item Für ERA hauptsächlich interessant:\\Binär ($B=2$), Oktal ($B=8=2^3$) und Hexadezimal ($B=16=2^4$)
		\item Was passiert bei einem Verschieben der Stellen nach links/rechts?\pause\\
		$\rightarrow$ nach links: $n\cdot B$, nach rechts: $\left\lfloor{\frac{n}{B}}\right\rfloor$
		\item Wie kann man Binärzahlen in Hexadezimal/Oktal umwandeln?\pause\\
		$\rightarrow$ 4 Binärziffern ergeben eine Hexadezimalziffer, 3 Binärziffern ergeben eine Oktalziffer
	\end{itemize}
\end{frame}

\begin{frame}[c]{Negative Zahlen}{}
	\begin{columns}[c]
		\begin{column}{0.3\textwidth}
			\begin{figure}
				\begin{tabular}{crr}
					&$0100\;1110_2$ & $=\phantom{-}78$ \\
					&\textcolor{TUMBlue}{$1$}$100\;1110_2$ & $=-78$
				\end{tabular}
				\caption{Vorzeichenbit}
			\end{figure}
		\end{column}
		\begin{column}{0.3\textwidth}
			\begin{figure}
				\begin{tabular}{crr}
					&$0100\;1110_2$ & $=\phantom{-}78$ \\
					&$1011\;0001_2$ &$=-78$
				\end{tabular}
				\caption{Einerkomplement}
			\end{figure}
		\end{column}
		\begin{column}{0.3\textwidth}
			\begin{figure}
				\begin{tabular}{rrl}
					&$0100\;1110_2$ & $=\phantom{-}78$ \\\hline
					&$1011\;0001_2$ &\\
					+ & $1_2$& \\\hline
					&$1011\;0010_2$ &$=-78$
				\end{tabular}
				\caption{Zweierkomplement}
			\end{figure}
		\end{column}
	\end{columns}
\end{frame}

\begin{frame}[c]{Rechnen im Binärsystem}{}
	\begin{itemize}
		\item schriftliches Rechnen grundsätzlich gleich wie im Dezimalsystem
		\item Subtraktion: Addition mit negativer Zahl (Zweierkomplement)
		\item Datenwortbreite bei negativen Zahlen ausschlaggebend!
		\item Ablesen von Zahlen im Zweierkomplement: Höchste Wertigkeit subtrahieren, Rest addieren
	\end{itemize}
	\vfill{}
	\begin{block}{Beispiel}
		\begin{enumerate}
			\item $(1100.0110)_2 = -2^7+(2^1+2^2+2^6)=-58$ bei Breite von 8 Bit
			\item $(1010)_2 = -2^3 + (2^1) = -6$ bei Breite von 4 Bit
		\end{enumerate}
	\end{block}
\end{frame}

\begin{frame}[c]{}{}
  \begin{center}
    \LARGE Fragen?
  \end{center}
\end{frame}

\begin{frame}[c]{Artemis-Hausaufgaben}{}
  \begin{center}
  	\LARGE kurze Einführung in Artemis/Git
  \end{center}
\end{frame}

\begin{frame}[fragile, c]{Links}{}
  \begin{itemize}
    \item Zulip: \href{https://zulip.cit.tum.de/#narrow/channel/3254-ERA-Tutorium-.E2.80.93-Mi-1600-1}{\enquote{ERA Tutorium -- Mi-1600-1}}
          bzw. \href{https://zulip.cit.tum.de/#narrow/channel/3266-ERA-Tutorium-.E2.80.93-Fr-1600-1}{\enquote{ERA Tutorium -- Fr-1600-1}}
    \item \href{https://www.moodle.tum.de/course/view.php?id=121343}{ERA-Moodle-Kurs}
    \item \href{https://artemis.tum.de/courses/596/exercises/21744}{ERA-Artemis-Kurs}
    \item \href{https://git-scm.com/docs/gittutorial}{Git-Tutorial}, \href{https://rogerdudler.github.io/git-guide/}{alternatives Tutorial}
  \end{itemize}
\end{frame}

\maketitle

\end{document}
